\definemathmatrix[pmatrix]        % defining matrix with parentheses
	[matrix:parentheses]
	[simplecommand=pmatrix]

\setvariables[meta]
  [title={Abstand: Punkt zu Gerade},
   subtitle={Anleitung},
   date={2025-03-20},
   author={Niklas von Hirschfeld},
  ]
\setupinteraction[title={\getvariable{meta}{title}}]

\starttext

\startfrontmatter
\component c_titlepage
% \component c_toc
\stopfrontmatter

\setuppagenumber[number=1]

\startbodymatter

\chapter{Abstand Punkt zu Ebene}

Um den Abstand von einem Punkt zu einer Ebene zu ermitteln, muss zunächst ein
Punkt auf der Ebene gefunden werden, welcher am {\em nächsten} an dem Punkt
ist. Von dort aus wird der Abstand zwischen den beiden Punkten berechnet
werden.

Um den Punkt auf der Ebene zu ermitteln geht man wie folgt vor:

\startitemize[packed,m]
\item Normalenvektor der Ebene bestimmen
\item Hilfsgerade mit Normalenvektor und Ortsvektor unseres gegeben Punkt aufstellen
\item Schnittpunkt zwischen der Gerade und der gegebenen Ebenen ermitteln
\stopitemize

\blank
\hrule
\blank

Je nach dem, welche Ebenenform gegeben ist, gibt es verschiedenen Wege an einen Normalenvektor heran zu kommen.

\startitemize[packed,m]
\item Bei der {\bf Parameterform} muss man das {\bf Keruzprodukt} der beiden {\bf Spannvektoren} berechnen. Daraus ergibt sich ein Normalenvektor. Wenn also folgende Ebenenform gegeben ist
\startformula
E: \vec{x} = \vec{a} + r \cdot \vec{b} + s \vec{c}
\stopformula
wird ein Normalenvektor wie folgt berechnet:
\startformula
\vec{n}=\vec{b}\times \vec{c}
\stopformula
\item Wenn die {\bf Koordinatenform} gegeben ist, lässt sich ein Normalenvektor aus dieser ablesen. Aus der Koordinatenform:
\startformula
E: n_1\cdot x_1 + n_2 \cdot x_2 + n_3 \cdot x_3 = d
\stopformula
geht folgender Normalenvektor hervor:
\startformula
\vec{n}=
\startpmatrix
\NC n_1 \NR 
\NC n_2 \NR
\NC n_3 \NR
\stoppmatrix
\stopformula
\item Wenn eine {\bf Normalenform} gegeben ist, kann man diesen aus der Form übernehmen.
\stopitemize

Mit hilfe eines Normalenvektor wird nun die Hilfgerade aufgestellt. Mit dem gegebene Puntk $P$ ist eine mögliche Hilfgerade folgende:

\startformula
g: \vec{x}=\vec{OP} + r \cdot \vec{n}
\stopformula

Diese Gerade hat die eigenschaften, das sie {\bf Orthogonal}\sidenote{{\bf
Orthogonal} bedeutet, dass der Vektor \quotation{senkrecht} zu dem andere
steht} durch unsere Ebene verläuft und durch unseren Punkt $P$. Das ist wichtig
für den folgenden Schritt. Da die Gerade die vorher genannten Eigentschaften
erfüllt, ist ihr Schnittpunkt mit der Ebene genau der Punkt mit dem geringsten
Abstand zu unserem Punkt $P$.

\startformula
g = E
\stopformula

Die dabei berechneten Parameter können in die Geraden- oder Ebenengleichung
eingesetzt werden, um den Schnittpunkt $S$ zu berechnen. Mit diesem Punkt
können wir den Abstand der Punkte $P$ und $S$ berechnen.

\startformula
|\overline{PS}|
\stopformula

\section{Abstand Punkt zu Punkt}

Der Abstand von einem Vektor ist wie folgt definiert:

\startformula
|\vec{a}|=\sqrt{a_1^2+a_2^2+a_3^2}
\stopformula

Um also den Abstand zwischen zwei Punkten $A$ und $B$ zu berechnen, muss der
Vektor $\vec{AB}$ zunächst ermittelt werden.
\startformula
\vec{AB}=\vec{OB} - \vec{OA}
\stopformula

\stopbodymatter

\startbackmatter
% \component c_definitions
% \component c_bibliography
\stopbackmatter

\stoptext

